\documentclass[10pt,landscape]{article}
\usepackage[landscape,margin=0.3in]{geometry}
\usepackage{multicol}
\usepackage{amsmath}
\usepackage{amssymb}
\usepackage{enumitem}
\usepackage{xcolor}
\usepackage[compact]{titlesec}

% Reduce spacing
\setlength{\parindent}{0pt}
\setlength{\parskip}{1pt}
\setlength{\columnsep}{10pt}

% Section formatting
\titleformat{\section}{\normalsize\bfseries\color{blue!70!black}}{\thesection}{0.5em}{}
\titlespacing*{\section}{0pt}{3pt}{1pt}
\titleformat{\subsection}{\small\bfseries\color{blue!50!black}}{\thesubsection}{0.5em}{}
\titlespacing*{\subsection}{0pt}{2pt}{1pt}

\begin{document}
\begin{multicols}{3}
\small

\begin{center}
\textbf{\large ECO310 Midterm Cheat Sheet - Clear Version}
\end{center}

\section{Core Utility Functions (Know These Cold!)}

\subsection{Cobb-Douglas: $u = x_1^a x_2^b$}
\textbf{Why it matters:} Appears on EVERY exam

\textbf{The "Magic Rule":} Spend a fraction of your income on each good based on the exponent ratios.
\begin{itemize}[leftmargin=*,noitemsep,topsep=0pt]
\item Fraction spent on good 1: $\frac{a}{a+b}$
\item Fraction spent on good 2: $\frac{b}{a+b}$
\end{itemize}

\textbf{Demand functions:}
$$x_1^* = \frac{a}{a+b} \cdot \frac{m}{p_1}, \quad x_2^* = \frac{b}{a+b} \cdot \frac{m}{p_2}$$

\textbf{Example:} $u = x_1^{0.2} x_2^{0.4}$, $m=15$, $p_1=1$, $p_2=2$: Normalize $0.2/0.6 = 1/3$. Spend $\$5$ on good 1 (5 units), $\$10$ on good 2 (5 units).

\textbf{MRS:} $MRS = -\frac{a}{b} \cdot \frac{x_2}{x_1}$ (slopes of indifference curves)

\textbf{At optimum:} $MRS = -\frac{p_1}{p_2}$ (slope of budget line)

\subsection{Perfect Complements: $u = \min\{ax_1, bx_2\}$}
\textbf{Why it matters:} Appears on EVERY exam

\textbf{Key insight:} You always consume at the "kink" where $ax_1 = bx_2$ (no point wasting money on extra of one good when you need them together).

\textbf{Demand functions:} Set $ax_1 = bx_2$ and use budget $p_1x_1 + p_2x_2 = m$:
$$x_1^* = \frac{m}{ap_2 + bp_1}, \quad x_2^* = \frac{am}{ap_2 + bp_1}$$

\textbf{Example:} $u = \min\{4x_1, x_2\}$, $m=20$, $p_1=p_2=1$: Set $4x_1 = x_2$, budget gives $x_1 + 4x_1 = 20 \Rightarrow x_1 = 4$, $x_2 = 16$.

\textbf{Critical fact for exams:} Perfect complements have \textbf{ZERO substitution effect} (entire price change is income effect, because you can't substitute between the goods).

\subsection{Perfect Substitutes: $u = ax_1 + bx_2$}
\textbf{Why it matters:} Tested in 3 out of 4 exams

\textbf{Key insight:} MRS is constant $= -a/b$. You only buy the good with better "bang-per-buck".

\textbf{Decision rule:} Compare $\frac{a}{p_1}$ (utility per dollar of good 1) vs $\frac{b}{p_2}$ (utility per dollar of good 2):
\begin{itemize}[leftmargin=*,noitemsep,topsep=0pt]
\item If $\frac{a}{p_1} > \frac{b}{p_2}$: Buy only good 1, $x_1^* = m/p_1$, $x_2^* = 0$
\item If $\frac{a}{p_1} < \frac{b}{p_2}$: Buy only good 2, $x_1^* = 0$, $x_2^* = m/p_2$
\item If equal: Any combination on budget line works
\end{itemize}

\textbf{Corner solutions:} This ALWAYS gives corner solutions unless bang-per-buck is exactly equal.

\subsection{CES Utility: $u = (x_1^r + x_2^r)^{1/r}$}
\textbf{Why it matters:} Tested in 2025 (and nests all three above!)

\textbf{Key insight:} As parameter changes, CES becomes:
\begin{itemize}[leftmargin=*,noitemsep,topsep=0pt]
\item $r \to -\infty$: Perfect complements
\item $r \to 0$: Cobb-Douglas
\item $r = 1$: Perfect substitutes
\end{itemize}

\textbf{MRS:} $MRS = -\left(\frac{x_2}{x_1}\right)^{1-r}$

\section{Finding Optimal Bundles (The Process)}

\subsection{Step 1: Set Up the Problem}
\textbf{Budget constraint:} $p_1 x_1 + p_2 x_2 = m$
\begin{itemize}[leftmargin=*,noitemsep,topsep=0pt]
\item Slope: $-p_1/p_2$
\item Intercepts: $(m/p_1, 0)$ and $(0, m/p_2)$
\end{itemize}

\subsection{Step 2: Find MRS}
\textbf{MRS (Marginal Rate of Substitution):} The rate at which you're willing to trade good 2 for good 1 while staying on the same indifference curve.

$$MRS = -\frac{dx_2}{dx_1}\Big|_{\text{same utility}} = -\frac{MU_1}{MU_2} = -\frac{\partial u/\partial x_1}{\partial u/\partial x_2}$$

\textbf{Interpretation:} MRS tells you your "personal price ratio" (how you value the goods). At optimum, this must equal the market price ratio.

\subsection{Step 3: Optimize (Interior Solution)}
\textbf{Tangency condition:} $MRS = -\frac{p_1}{p_2}$

Or equivalently: $\frac{MU_1}{MU_2} = \frac{p_1}{p_2}$ or $\frac{MU_1}{p_1} = \frac{MU_2}{p_2}$

\textbf{What this means:} "Bang-per-buck" must be equal across all goods. If good 1 gave you more utility per dollar, you'd buy more of it.

\textbf{Budget binds:} $p_1 x_1 + p_2 x_2 = m$ (spend all your money)

\textbf{Solve these two equations together} to get $x_1^*$ and $x_2^*$.

\subsection{Step 4: Check for Corner Solutions}
\textbf{When do corners happen?}
\begin{itemize}[leftmargin=*,noitemsep,topsep=0pt]
\item Perfect substitutes (always a corner unless bang-per-buck exactly equal)
\item When your FOC (tangency condition) gives you a negative quantity
\item When goods aren't worth the price (MRS never equals price ratio)
\end{itemize}

\textbf{How to check:} If tangency gives $x_i < 0$, set $x_i = 0$ and spend all $m$ on the other good.

\section{Marshallian vs Hicksian Demand}

\subsection{Marshallian Demand: $x_i(p, m)$}
\textbf{What is it:} The regular demand function. What you buy given prices and income.

\textbf{Problem it solves:} Maximize utility $u(x_1, x_2)$ subject to budget $p_1x_1 + p_2x_2 = m$

\textbf{Result:} Gives you the best bundle you can afford.

\subsection{Hicksian (Compensated) Demand: $h_i(p, \bar{u})$}
\textbf{What is it:} The demand that achieves a target utility level at minimum cost.

\textbf{Problem it solves:} Minimize expenditure $p_1x_1 + p_2x_2$ subject to utility constraint $u(x_1, x_2) = \bar{u}$

\textbf{Result:} Gives you the cheapest bundle that achieves utility $\bar{u}$.

\textbf{When you use it:} For substitution effects (holding utility constant while prices change).

\subsection{Example: Finding Hicksian Demand}
$u = x_1 x_2$, target $\bar{u} = 64$, $p_1 = p_2 = 4$:

MRS = price ratio: $x_2/x_1 = 1 \Rightarrow x_1 = x_2$

Use $u = 64$: $x_1^2 = 64 \Rightarrow x_1 = x_2 = 8$

\section{Substitution \& Income Effects}

\subsection{The Big Picture}
When price of good 1 rises, two things happen:
\begin{enumerate}[leftmargin=*,noitemsep,topsep=0pt]
\item \textbf{Substitution effect:} Good 1 is now relatively more expensive, so you substitute away from it (even keeping utility constant)
\item \textbf{Income effect:} Your purchasing power decreased, so you adjust consumption based on whether good 1 is normal or inferior
\end{enumerate}

\textbf{Total effect = Substitution effect + Income effect}

\subsection{The Recipe (Always Follow These Steps!)}
Suppose $p_1$ increases from $p_1^{\text{old}}$ to $p_1^{\text{new}}$:

\textbf{Step 1:} Find \textbf{original bundle} $(x_1^O, x_2^O)$ at old prices

\textbf{Step 2:} Find \textbf{new bundle} $(x_1^N, x_2^N)$ at new prices

\textbf{Step 3:} Find \textbf{compensated bundle} $(x_1^C, x_2^C)$:
\begin{itemize}[leftmargin=*,noitemsep,topsep=0pt]
\item Use new prices $p_1^{\text{new}}$, $p_2$
\item Keep utility at original level: $u(x_1^C, x_2^C) = u(x_1^O, x_2^O)$
\item This is the Hicksian demand at new prices!
\end{itemize}

\textbf{Step 4:} Calculate effects:
\begin{itemize}[leftmargin=*,noitemsep,topsep=0pt]
\item Substitution effect: $\Delta x_1^{\text{sub}} = x_1^C - x_1^O$
\item Income effect: $\Delta x_1^{\text{inc}} = x_1^N - x_1^C$
\item Total effect: $\Delta x_1^{\text{total}} = x_1^N - x_1^O$
\end{itemize}

\subsection{Sign Rules}
\textbf{Substitution effect:} ALWAYS negative (or zero) when price rises
\begin{itemize}[leftmargin=*,noitemsep,topsep=0pt]
\item If $p_1$ rises, $x_1^C < x_1^O$ (you buy less to substitute away)
\item Exception: Perfect complements have zero substitution effect
\end{itemize}

\textbf{Income effect:} Depends on type of good
\begin{itemize}[leftmargin=*,noitemsep,topsep=0pt]
\item Normal good: Negative (buy less when poorer)
\item Inferior good: Positive (buy more when poorer)
\end{itemize}

\subsection{Example with Cobb-Douglas}
$u = BT$, $m=64$, $p_B$ rises from $\$1$ to $\$4$, $p_T = \$4$

\textbf{Original:} $B^O = 32$, $T^O = 8$, $u = 256$

\textbf{New:} $B^N = 8$, $T^N = 8$

\textbf{Compensated:} $B = T$ and $BT = 256 \Rightarrow B^C = 16$, $T^C = 16$

\textbf{Effects:} Sub: $-16$, Income: $-8$, Total: $-24$

\section{Expected Utility (Key for Every Exam!)}

\subsection{Setup}
\textbf{Lottery:} Get outcome $c_i$ with probability $p_i$

\textbf{Expected utility:}
$$EU(L) = \sum_{i} p_i \cdot u(c_i) = p_1 u(c_1) + p_2 u(c_2) + \ldots$$

\textbf{Critical:} Use final wealth in each state, not the change!

\subsection{Example: Setting Up EU}
Have $\$100$, invest $\$x$ in risky asset (+100\% or -50\%, each prob 1/2):

Good: $W = 100 + x$, Bad: $W = 100 - 0.5x$

$EU = \frac{1}{2}u(100+x) + \frac{1}{2}u(100-0.5x)$

\subsection{Optimizing EU}
\textbf{To find optimal $x$:}
\begin{enumerate}[leftmargin=*,noitemsep,topsep=0pt]
\item Take derivative: $\frac{d(EU)}{dx}$
\item Set equal to zero: $\frac{d(EU)}{dx} = 0$
\item Solve for $x^*$
\item Check corners: $x = 0$ or $x = $ max possible
\end{enumerate}

\subsection{Risk Aversion}
\textbf{Definition:} You prefer certainty to a gamble with the same expected value.

\textbf{Math version:} $u(E[W]) > EU$ (utility of expected wealth $>$ expected utility)

This happens when $u$ is \textbf{concave} ($u'' < 0$)

\textbf{Risk neutral:} $u$ is linear, $u(E[W]) = EU$

\textbf{Risk loving:} $u$ is convex ($u'' > 0$), $u(E[W]) < EU$

\subsection{Certainty Equivalent (CE)}
\textbf{Definition:} The guaranteed amount that gives the same utility as the lottery.

$$u(CE) = EU$$

\textbf{Risk premium:} $RP = E[W] - CE$
\begin{itemize}[leftmargin=*,noitemsep,topsep=0pt]
\item Risk averse: $RP > 0$ (would pay to avoid risk)
\item Risk neutral: $RP = 0$
\end{itemize}

\textbf{Example:} 50-50 of $\$0$ or $\$200$, $u(w) = \sqrt{w}$: $EU = \frac{1}{2}\sqrt{200}$, so $CE = 50$. $E[W] = 100$, thus $RP = 50$.

\subsection{Common Utility Functions}
\textbf{Log utility:} $u(w) = \ln(w)$
\begin{itemize}[leftmargin=*,noitemsep,topsep=0pt]
\item Risk averse (concave)
\item Used frequently on exams
\end{itemize}

\textbf{CRRA:} $u(w) = \frac{w^{1-\rho}}{1-\rho}$ (or $\ln w$ if $\rho=1$)
\begin{itemize}[leftmargin=*,noitemsep,topsep=0pt]
\item $\rho > 0$: Risk averse
\item Higher $\rho$: More risk averse
\end{itemize}

\section{WARP (Revealed Preference)}

\subsection{What is WARP?}
\textbf{Weak Axiom of Revealed Preference:} If you choose bundle $x$ when bundle $y$ is affordable, then whenever you choose $y$, bundle $x$ must NOT be affordable (or must cost more).

\textbf{Interpretation:} Your choices should be consistent. If you "revealed" you prefer $x$ to $y$ by choosing it when both were available, you shouldn't later choose $y$ when $x$ is also available.

\subsection{How to Check WARP (Two Periods)}
Given: Period A bundle $(x_1^A, x_2^A)$, prices $(p_1^A, p_2^A)$, spent $m^A$

Period B bundle $(x_1^B, x_2^B)$, prices $(p_1^B, p_2^B)$, spent $m^B$

\textbf{Step 1:} Check if A's bundle was revealed preferred to B's bundle:
\begin{itemize}[leftmargin=*,noitemsep,topsep=0pt]
\item Cost of B's bundle at A's prices: $p_1^A x_1^B + p_2^A x_2^B$
\item If this $\leq m^A$, then A's choice revealed preference for A's bundle
\end{itemize}

\textbf{Step 2:} Check if B's bundle was revealed preferred to A's bundle:
\begin{itemize}[leftmargin=*,noitemsep,topsep=0pt]
\item Cost of A's bundle at B's prices: $p_1^B x_1^A + p_2^B x_2^A$
\item If this $\leq m^B$, then B's choice revealed preference for B's bundle
\end{itemize}

\textbf{WARP violated if:} Both bundles are revealed preferred (at least one strictly)

\subsection{Example}
Mon: Spent $\$90$, Tue bundle cost $\$80$ at Mon prices $\Rightarrow$ Mon strictly revealed preferred

Tue: Spent $\$100$, Mon bundle cost $\$100$ at Tue prices $\Rightarrow$ Tue weakly revealed preferred

\textbf{Result:} WARP violated (both revealed preferred)

\section{Corner Solutions \& Non-Negativity}

\subsection{When Do Corners Happen?}
\textbf{Situation 1:} Perfect substitutes (almost always a corner)

\textbf{Situation 2:} When FOC gives negative quantity
\begin{itemize}[leftmargin=*,noitemsep,topsep=0pt]
\item You set MRS = price ratio, solve, get $x_i < 0$
\item This means at any positive quantity, the good isn't worth buying
\item Solution: Set $x_i = 0$, spend all income on other good
\end{itemize}

\textbf{Situation 3:} MRS never equals price ratio
\begin{itemize}[leftmargin=*,noitemsep,topsep=0pt]
\item If $|MRS| > p_1/p_2$ everywhere: consume only good 1
\item If $|MRS| < p_1/p_2$ everywhere: consume only good 2
\end{itemize}

\subsection{How to Check in Your Answer}
\textbf{ALWAYS verify:} $x_1 \geq 0$ and $x_2 \geq 0$

If either is negative, you made an error or you're at a corner.

\section{Quick Tips for Common Exam Traps}

\begin{enumerate}[leftmargin=*,noitemsep,topsep=0pt]
\item \textbf{Cobb-Douglas:} "Spend 3x as much" $\neq$ "Buy 3x as much" (unless prices are equal)

\item \textbf{Perfect complements:} Substitution effect is ALWAYS zero (entire change is income effect)

\item \textbf{Portfolio problems:} Variance $\sigma^2 \neq$ Standard deviation $\sigma$

\item \textbf{EU setup:} Use FINAL wealth in each state, not the change

\item \textbf{Marshallian vs Hicksian:} Marshallian maximizes utility given income; Hicksian minimizes cost given utility

\item \textbf{Risk aversion:} Check if $u(E[W]) > EU$, not just if $u$ is concave

\item \textbf{WARP:} "Revealed preferred" = chosen when other was affordable (cost $\leq$ budget)

\item \textbf{Substitution effect:} ALWAYS use compensated bundle (same utility, new prices)
\end{enumerate}

\end{multicols}
\end{document}
